\documentclass{article}

\usepackage[utf8]{inputenc}
\usepackage{amsmath}
\usepackage{natbib}

\title{Counterexample-Guided Formal Verification for Agentic Text-to-SQL Refinement}

\author{Anonymous}

\begin{document}

\maketitle

\begin{abstract}
LLM-based Text-to-SQL agents solve database query tasks through an agentic loop of
retrieval, candidate generation, verification, and refinement. Yet refining an
erroneous SQL query without introducing new correctness bugs remains an open
challenge, and the termination conditions used by prior systems---mere execution
success \cite{wang2023macsql} or LLM-generated unit tests \cite{talaei2024chess}---are
necessary but not sufficient to guarantee a correct query. Formal verification has
been applied to Text-to-SQL only for \emph{offline evaluation}, where SpotIt checks
bounded equivalence of a generated query against a human gold query
\cite{klopfenstein2026spotit}; no prior method uses it to guide refinement online,
where no gold query exists. We propose to close this gap with a
formal-verification-guided refinement scheme: a \textbf{Critic} agent first states
the correctness criteria a query must satisfy given the task and schema, and a
bounded SMT-based verifier searches for a counterexample database that
differentiates the candidate from those criteria. Any such counterexample is
returned to the generator as concrete repair evidence, and verification---rather than
execution success---terminates the loop. The central difficulty, and our primary
research lever, is the \emph{specification gap}: soundness holds only relative to the
correctness criteria, so our contribution is a tractable way to synthesize
verifiable criteria from the question and schema---distinct from full SQL
generation---together with counterexample-guided repair. In contrast to
heuristic error detection \cite{hong26errorllm}, our verifier yields a sound
differentiating witness. We plan to evaluate on the BIRD benchmark
\cite{li2023bird} and hypothesize that counterexample-guided verification improves
execution accuracy and, more importantly, closes the silent-error gap that
test-based evaluation overlooks \cite{klopfenstein2026spotit}.
\end{abstract}

\bibliographystyle{plain}
\bibliography{ref}

\end{document}
